% Part: sets-functions-relations
% Chapter: arithmetization
% Section: From N to Z

\documentclass[../../../include/open-logic-section]{subfiles}


\begin{document}

\olfileid{sfr}{arith}{int}
\olsection{Van $\Nat$ naar $\Int$}
We beginnen met twee eenvoudige observaties:
	\begin{enumerate}
		\item Elk geheel getal kan worden geschreven als $n - m$, waarbij $n, m \in\Nat$.
		\item De informatie in een uitdrukking $n - m$ kan evengoed worden vastgelegd met een geordend paar $\tuple{n, m}$.
	\end{enumerate}
We weten al dat de geordende paren van natuurlijke getallen de !!{element}s
van $\Nat^2$ zijn. Bovendien nemen we aan dat we $\Nat$ begrijpen. De
twee observaties leiden daarom tot een na\"{i}ef voorstel: \emph{vat
gehele getallen op als geordende paren van natuurlijke getallen}.
Dit voorstel is echter te na\"{i}ef. We willen uiteraard dat $0- 2 = 4
- 6$, maar $\tuple{0, 2 }\neq \tuple{4, 6}$. We kunnen $\Nat^2$ dus
niet zonder meer de verzameling gehele getallen noemen.
Als we het probleem uit dit voorbeeld veralgemenen, willen we het volgende:
	$$a - b = c - d \text{ precies dan als }a + d = c + b$$
(Dat dit de \emph{bedoelde} werking is, ligt voor de hand: tel aan beide
kanten $b$ en $d$ op.) De eenvoudigste manier om dit gedrag te
garanderen is de equivalentierelatie $\Intequiv$ tussen geordende paren
als volgt te definiëren:
$$\tuple{ a, b } \Intequiv \tuple{c, d}\text{ precies dan als }a + d = c + b$$
We moeten nu aantonen dat dit een equivalentierelatie is.
\begin{prop} $\Intequiv$ is een equivalentierelatie.
\end{prop}
	\begin{proof}
	We moeten aantonen dat $\Intequiv$ reflexief, symmetrisch en transitief is.
	\emph{Reflexiviteit:} Het is duidelijk dat $\tuple{a, b} \Intequiv \tuple{a, b}$,
	want $a + b = b + a$.
	\emph{Symmetrie:} Stel dat $\tuple{a, b} \Intequiv \tuple{c, d}$,
	zodat $a + d = c + b$. Dan is $c + b = a + d$, en dus
	$\tuple{c, d} \Intequiv \tuple{a, b}$.
	\emph{Transitiviteit:} Stel dat $\tuple{a, b} \Intequiv \tuple{c,
	d}\Intequiv \tuple{m, n}$. Dan is $a + d = c + b$ en $c + n = m +
	d$. Bijgevolg is $a + d + c + n = c + b + m + d$, en dus $a + n =
	m + b$. Hieruit volgt $\tuple{a, b } \Intequiv \tuple{m, n}$.
\end{proof}
Met deze equivalentierelatie kunnen we nu equivalentieklassen vormen:
\begin{defn}
		De gehele getallen zijn de equivalentieklassen van geordende paren
		van natuurlijke getallen onder $\Intequiv$; dus $\Int =
		\equivclass{\Nat^2}{\Intequiv}$.
\end{defn}
Men kan op uiteenlopende manieren \emph{filosofisch} reageren op deze
stipulatieve definitie. Voordat we die reacties bespreken, is het echter
nuttig eerst enkele technische details uit te werken.
Nu we hebben bepaald wat de gehele getallen zijn, moeten we de elementaire
functies en relaties daarop definiëren. We schrijven
$\equivrep{m, n}{\Intequiv}$ voor de equivalentieklasse onder
$\Intequiv$ die $\tuple{m, n}$ als !!a{element} bevat.\footnote{Met de
in \olref[sfr][rel][eqv]{def:equivalenceclass} ingevoerde notatie
zouden we hiervoor $\equivrep{\tuple{m,n}}{\Intequiv}$ hebben
geschreven. Dat is echter iets moeilijker leesbaar.} Dus:
	$$\equivrep{m, n}{\Intequiv} = \Setabs{\tuple{a, b} \in \Nat^2}{\tuple{a, b}\Intequiv \tuple{m, n}}$$
We definiëren nu:
	\begin{align*}
		\equivrep{a, b}{\Intequiv} + \equivrep{c, d}{\Intequiv} &= \equivrep{a + c, b + d}{\Intequiv}\\
		\equivrep{a, b}{\Intequiv} \times \equivrep{c, d}{\Intequiv} &= \equivrep{a c + b  d, a  d + b c}{\Intequiv}\\
		\equivrep{a, b}{\Intequiv} \leq \equivrep{c, d}{\Intequiv} &\text{ precies dan als }a + d \leq b + c
	\end{align*}	
(Zoals gebruikelijk schrijf ik `$ab$' voor `$(a \times b)$', zodat de
axioma's gemakkelijker leesbaar zijn.) We moeten nu nagaan of deze
definities naar behoren \emph{werken}. Het
is tamelijk omslachtig om dit volledig te formuleren en te controleren;
de details staan in \olref[check]{sec}. De kern is dat alles werkt.
Er rest nog één punt. We hebben de gehele getallen met behulp van
natuurlijke getallen geconstrueerd. Dit betekent echter dat de
natuurlijke getallen \emph{zelf geen gehele getallen zijn}. Op de
filosofische betekenis hiervan komen we terug in \olref[ref]{sec}.
Technisch hebben we een manier nodig om natuurlijke getallen
\emph{als} gehele getallen te behandelen. De oplossing is eenvoudig:
voor elke $n \in \Nat$ spreken we af dat $n_\Int =
\equivrep{n, 0}{\Intequiv}$. We
moeten controleren of deze afspraak goed werkt, dat wil
zeggen dat voor alle $m, n \in \Nat$ geldt:
\begin{align*}
	(m + n)_\Int &= m_\Int + n_\Int\\
	(m \times n)_\Int &= m_\Int \times n_\Int\\
	m \leq n &\liff m_\Int \leq n_\Int
\end{align*}
Dit is vrij eenvoudig na te gaan. Voor de tweede gelijkheid
kunnen we bijvoorbeeld de eigenschappen van de natuurlijke getallen
gebruiken en als volgt redeneren:
%\begin{proof}
%	Met de eigenschappen van de natuurlijke getallen is duidelijk:
	\begin{align*}
%		(m + n)_\Int  = \equivrep{m + n, 0}{\Intequiv} = \equivrep{m + n, 0 + 0}{\Intequiv} = \equivrep{m, 0}{\Intequiv} + \equivrep{n, 0}{\Intequiv} = m_\Int + n_\Int\\
		(m \times n)_\Int  &= \equivrep{m \times n, 0}{\Intequiv} \\
		&= \equivrep{m \times n + 0 \times 0, m \times 0 + 0 \times n}{\Intequiv} \\
		&= \equivrep{m, 0}{\Intequiv} \times \equivrep{n, 0}{\Intequiv} \\
		&= m_\Int \times n_\Int
	\end{align*}
%\end{proof}
De controle van de overige twee voorwaarden laten we als oefening.
\begin{prob}
	Toon aan dat $(m + n)_\Int = m_\Int + n_\Int$ en $m \leq n \liff
	m_\Int \leq n_\Int$ voor alle $m, n \in \Nat$.
\end{prob}
\end{document}
